\section*{Capítulo \ref{ch:g12:metamorphic-rocks} --- As rochas metamórficas}

\begin{solution}{exo:g12:metamorphic-rocks:1}
(a) \omterm{def:g12:metamorphic-rocks:polymorph}{polimorfos}; (b) uma \omterm{def:g12:metamorphic-rocks:index-mineral}{isógrada}; (c) uma \omterm{def:g12:metamorphic-rocks:contact-regional}{auréola}; (d) uma \omterm{def:g12:metamorphic-rocks:foliation}{foliação};
(e) a \omterm{def:g12:metamorphic-rocks:lithostatic}{pressão litostática}.
\end{solution}

\begin{solution}{exo:g12:metamorphic-rocks:2}
(1) mármore; (2) quartzito; (3) ardósia; (4) gnaisse; (5) hornfels.
\end{solution}

\begin{solution}{exo:g12:metamorphic-rocks:3}
$P = \rho g z$: a \qty{10}{km}, $2830 \times 9.81 \times \num{10000} \approx
\qty{2.78e8}{Pa} = \qty{0.278}{GPa}$; a \qty{20}{km}, o dobro,
\qty{0.555}{GPa}.
\end{solution}

\begin{solution}{exo:g12:metamorphic-rocks:4}
(1) Cianita (a linha cianita--andaluzita fica perto de \qty{0.25}{GPa} a
\qty{400}{\celsius}). (2) Andaluzita. (3) Silimanita.
\end{solution}

\begin{solution}{exo:g12:metamorphic-rocks:5}
(1) Ardósia, filito, xisto, gnaisse. (2) No xisto. (3) O gnaisse.
\end{solution}

\begin{solution}{exo:g12:metamorphic-rocks:6}
(1) $2830 \times 9.81 \times \num{35000} \approx \qty{9.72e8}{Pa}$, isto é,
$\num{9.72e8}/101325 \approx \num{9600}$ atmosferas. (2) $1025 \times 9.81
\times 4000 \approx \qty{4.02e7}{Pa} = \qty{40.2}{MPa}$. (3) $972/40.2
\approx 24$ vezes.
\end{solution}

\begin{solution}{exo:g12:metamorphic-rocks:7}
(1) $z = \num{0.375e9}/(2830 \times 9.81) \approx \qty{13.5}{km}$.
(2) $(504 - 10)/13.5 \approx \qty{36.6}{\celsius}$ por quilômetro. (3) Não:
a rocha teria de ficar exatamente num ponto do diagrama; rochas assim são
raras.
\end{solution}

\begin{solution}{exo:g12:metamorphic-rocks:8}
(1) $23/3 \approx \qty{8}{mm}$. (2) Cerca de 8 milhões de anos. (3) A rocha
ficou em profundidade, a uma temperatura alta, durante milhões de anos.
\end{solution}

\begin{solution}{exo:g12:metamorphic-rocks:9}
(1) Quatro: as \omterm{def:g12:metamorphic-rocks:index-mineral}{isógradas} da granada, da estaurolita, da cianita e da
silimanita, nessa ordem. (2) Para oeste. (3) Na extremidade oeste, na zona
da silimanita. As suas rochas cruzaram a linha cianita--silimanita:
enterradas a cerca de \qty{20}{km} (\qty{0.56}{GPa}), elas estavam mais
quentes que cerca de \qty{590}{\celsius}.
\end{solution}

\begin{solution}{exo:g12:metamorphic-rocks:10}
(1) A lama vira hornfels, o calcário vira mármore, o arenito vira
quartzito. (2) Elas foram aquecidas sem ser comprimidas: nenhuma pressão
dirigida alinhou os seus minerais. (3) As rochas em volta dele já estão
quentes, então o magma quase não as aquece mais.
\end{solution}

\begin{solution}{exo:g12:metamorphic-rocks:11}
(1) Atravessadas ao empurrão: as suas faces planas estão voltadas para leste
e oeste, então as escamas ficam em planos norte--sul. (2) Vertical, orientada
norte--sul. (3) Os grãos de calcita e de quartzo não são nem achatados nem
alongados, então, mesmo quando se recristalizam sob compressão, não há nada
para alinhar.
\end{solution}

\begin{solution}{exo:g12:metamorphic-rocks:12}
(1) $z = (T - 10)/30$, logo $P = 0.0278\,(T - 10)/30 \approx
0.000927\,(T - 10)$. (2) $0.000927\,(T - 10) = 0.577 + 0.00209\,(T - 600)$
dá $0.001163\,T \approx 0.668$, logo $T \approx \qty{574}{\celsius}$, $P
\approx 0.000927 \times 564 \approx \qty{0.523}{GPa}$ e $z = (574 - 10)/30
\approx \qty{18.8}{km}$. (3) Cianita acima dessa profundidade, silimanita
abaixo dela.
\end{solution}

\begin{solution}{exo:g12:metamorphic-rocks:13}
(1) $10 + 10 \times 35 = \qty{360}{\celsius}$. (2) $10 + 30 \times 35 =
\qty{1060}{\celsius}$. (3) A pressão a \qty{35}{km} é de cerca de
\qty{0.97}{GPa}; a \qty{360}{\celsius}, esse ponto fica bem dentro do
campo da cianita. O segundo ponto fica muito à direita do diagrama, no
campo da silimanita, a temperaturas em que as rochas da crosta começam a se
fundir. (4) Eles precisam de uma pressão alta a uma temperatura baixa, uma
combinação que só a crosta fria levada para baixo numa \omterm{def:g11:subduction-zones:subduction}{zona de subducção}
oferece.
\end{solution}

\begin{solution}{exo:g12:metamorphic-rocks:14}
(1) Cianita $1000/3.61 \approx \qty{277}{cm^3}$; silimanita $1000/3.24
\approx \qty{309}{cm^3}$; andaluzita $1000/3.15 \approx \qty{317}{cm^3}$.
(2) $(317 - 277)/317 \approx 0.127$, cerca de \qty{13}{\percent}. (3) Sob alta
pressão, o arranjo mais compacto dos átomos é favorecido; a cianita
acomoda os mesmos átomos no menor espaço.
\end{solution}

\begin{solution}{exo:g12:metamorphic-rocks:15}
(1) \qty{0.555}{GPa} a \qty{20}{km}, \qty{0.694}{GPa} a \qty{25}{km}.
(2) $T = 600 + (P - 0.577)/0.00209$: cerca de \qty{590}{\celsius} a
\qty{0.555}{GPa} e \qty{656}{\celsius} a \qty{0.694}{GPa}. A
\qty{700}{\celsius}, as rochas estavam mais quentes que ambos, no campo da
silimanita: coerente. (3) $(700 - 10)/20 \approx \qty{34.5}{\celsius}$ e
$(700 - 10)/25 \approx \qty{27.6}{\celsius}$ por quilômetro, perto do
gradiente típico. (4) De $\qty{20}{km}$ a $\qty{25}{km}$ em 40 milhões de
anos: de 0.50 a \qty{0.63}{km} por milhão de anos, isto é, de 0.50 a
\qty{0.63}{mm} por ano.
\end{solution}

\begin{solution}{pb:g12:metamorphic-rocks:1}
\textbf{1.} Ela é foliada, e as suas micas são grandes o bastante para
serem vistas. As micas, a muscovita e a biotita, formam a \omterm{def:g12:metamorphic-rocks:foliation}{xistosidade}.

\textbf{2.} \ce{Al2SiO5}; são \omterm{def:g12:metamorphic-rocks:polymorph}{polimorfos}.

\textbf{3.} Os volumes de um quilograma estão na razão $3.15/3.61 \approx
0.873$: a cianita é cerca de \qty{13}{\percent} menor. A pressão favorece a
cianita.

\textbf{4.} A zona da cianita. Na descida, ela cruzou as \omterm{def:g12:metamorphic-rocks:index-mineral}{isógradas} da
biotita, da granada, da estaurolita e da cianita.

\textbf{5.} Cerca de 8 milhões de anos.

\textbf{6.} $P = 0.577 + 0.00209 \times 0 = \qty{0.577}{GPa}$.

\textbf{7.} \qty{0.70}{GPa} fica acima de \qty{0.577}{GPa}, do lado da
cianita da linha cianita--silimanita; fica também muito acima do campo da
andaluzita, que a \qty{600}{\celsius} fica abaixo de cerca de
\qty{0.23}{GPa}.

\textbf{8.} \qty{0.577}{GPa}.

\textbf{9.} $0.70 = 0.577 + 0.00209\,(T - 600)$, logo $T = 600 +
0.123/0.00209 \approx \qty{659}{\celsius}$.

\textbf{10.} \qty{7.0}{kbar}; $\num{7.0e8}/101325 \approx \num{6900}$
atmosferas.

\textbf{11.} $\rho g = 2830 \times 9.81 \approx \qty{27800}{Pa}$ por metro,
\qty{27.8}{MPa} por quilômetro.

\textbf{12.} $z = \num{0.70e9}/\num{27762} \approx \qty{25200}{m}$, cerca de
\qty{25.2}{km}.

\textbf{13.} $\num{0.577e9}/\num{27762} \approx \qty{20.8}{km}$.

\textbf{14.} $\num{0.70e9}/(2700 \times 9.81) \approx \qty{26.4}{km}$, cerca
de \qty{5}{\percent} mais fundo ($26.4/25.2 \approx 1.048$).

\textbf{15.} Na crosta: \qty{25}{km} é menos que até mesmo uma espessura
normal da crosta.

\textbf{16.} $(600 - 10)/25.2 \approx \qty{23.4}{\celsius}$ por quilômetro.

\textbf{17.} Ele é muito mais quente que uma \omterm{def:g11:subduction-zones:subduction}{zona de subducção} (abaixo de
\qty{10}{\celsius} por quilômetro) e um pouco mais frio que uma crosta
típica (\qty{30}{\celsius}), longe de uma região vulcânica (de 40 a
\qty{50}{\celsius}): uma rocha enterrada fundo na crosta espessada de um
cinturão de montanhas, por \omterm{def:g12:metamorphic-rocks:contact-regional}{metamorfismo regional}.

\textbf{18.} $10 + 30 \times 25.2 \approx \qty{766}{\celsius}$. A linha
cianita--silimanita, prolongada, fica a $0.577 + 0.00209 \times 166
\approx \qty{0.92}{GPa}$ a essa temperatura; \qty{0.70}{GPa} fica abaixo
dela, no campo da silimanita. Uma \omterm{def:g11:lithosphere-asthenosphere:geotherm}{geoterma} típica teria produzido
silimanita.

\textbf{19.} $\qty{25.2}{km} = \qty{2.52e7}{mm}$; a \qty{0.5}{mm} por ano,
cerca de 50 milhões de anos.

\textbf{20.} À medida que a rocha subia, ela esfriava, então as suas reações
ficaram mais lentas, e a água de que as reações inversas precisam tinha sido
expulsa na descida.

\textbf{21.} O xisto com cianita se cristalizou a cerca de
\textbf{\qty{25}{km} de profundidade}, a cerca de \qty{600}{\celsius}.
\end{solution}
